% !TeX program = pdflatex
% papers/conversa.tex — Corpus casual (baixa complexidade).
% Teoria fica para depois: primeiro a assistente aprende a falar educadamente.
% Ingerido por tools/ingere.py + tools/casual.sh; gabarito banal: tools/gabarito_llama.sh
\documentclass[11pt,a4paper]{article}
\usepackage[utf8]{inputenc}\usepackage[T1]{fontenc}\usepackage[brazil]{babel}
\usepackage[margin=2.4cm]{geometry}
\usepackage[hidelinks]{hyperref}
\newcommand{\code}[1]{\texttt{\small #1}}
% ─────────────────────────────────────────────────────────────────────────────
%  papers/gkcapa.tex — a capa do Reino, mínima, para os papers em `article`.
%
%  O estilo.tex completo é do `report` (títulos de capítulo, skak, …). Aqui fica
%  só o que a capa precisa, para o paper continuar a compilar sozinho com
%  pdflatex. O tradutor próprio (tests/tex) carrega o estilo.tex da raiz / o
%  symlink papers/estilo.tex e avalia o mesmo `\gkcapa`.
% ─────────────────────────────────────────────────────────────────────────────
\definecolor{ouro}{HTML}{B8912F}
\definecolor{ouroclaro}{HTML}{F3E9CE}
\definecolor{tinta}{HTML}{1E1B16}
\definecolor{regua}{HTML}{6E6858}
\providecommand{\gknota}{\fontsize{7.62}{11.05}\selectfont}
\providecommand{\gktexto}{\fontsize{10.50}{15.22}\selectfont}
\providecommand{\gksub}{\fontsize{12.33}{17.87}\selectfont}
\providecommand{\gksec}{\fontsize{14.47}{20.98}\selectfont}
\providecommand{\gkcap}{\fontsize{16.99}{24.63}\selectfont}
\providecommand{\gktit}{\fontsize{23.42}{33.95}\selectfont}
\providecommand{\Prj}{\mathbb{P}}
\providecommand{\Trf}{\mathcal{T}}
\providecommand{\gkcapa}[3]{%
\title{\vspace{-0.6cm}
{\color{ouro}\rule{\textwidth}{1.4pt}}\\[6mm]
\textbf{\gktit\color{tinta} Reino Dourado}\\[3mm]{\gksec\itshape\color{regua} Golden Kingdom}\\[8mm]
{\gkcap\scshape\color{ouro} #1}\\[5mm]
{\gksub\color{regua} #2}\\[2mm]
{\gktexto\color{regua}\itshape #3}\\[7mm]
{\color{ouro}\rule{\textwidth}{0.6pt}}\\[9mm]{\gknota\color{regua}\begin{minipage}{\textwidth}\setlength{\parindent}{0pt}%
\textbf{Aviso de Isenção Algorítmica:} O universo, as leis e as entidades contidas nestas páginas ---
incluindo felinos matriciais, aves de rapina algébricas e amuletos de controle de fluxo --- são
construções de pura ficção formal. Esta obra não descreve a realidade física, não serve como
engenharia reversa do cosmos e não constitui doutrina religiosa. Qualquer semelhança com bugs reais do seu
sistema operacional é mera coincidência (ou falta de reza). Prossiga por sua própria conta e
risco.\end{minipage}}}
\author{}\date{}}


\begin{document}
\gkcapa{A Conversa}
       {corpus casual --- assuntos banais}
       {educação no trato; complexidade sobe depois}
\maketitle

\begin{abstract}
\noindent Corpus de \textbf{conversa banal}: cumprimentos, rotina, comida, tempo, humor leve.
Sem teoria nesta fase. O gabarito (\texttt{tools/gabarito\_llama.sh}) enche a banda com pares
do dia a dia; a doutrina entra quando a assistente já souber falar.
\end{abstract}

\section{bom dia}
Bom dia! Espero que tenhas um dia leve. Em que posso ajudar?

\section{boa tarde}
Boa tarde. Tudo bem por aí?

\section{boa noite}
Boa noite. Descansa bem.

\section{oi}
Oi! Tudo bem?

\section{olá}
Olá. Prazer em falar contigo.

\section{oi tudo bem}
Oi, tudo bem sim --- e contigo?

\section{tudo bem}
Tudo bem, obrigado. E tu?

\section{como estás}
Estou bem, obrigado por perguntares. E tu, como vais?

\section{como vai}
Vou bem. Conta-me: como foi o teu dia?

\section{obrigado}
De nada. Estou por aqui.

\section{obrigada}
De nada. Foi um prazer.

\section{tchau}
Tchau! Cuida-te.

\section{até logo}
Até logo. Foi bom conversar.

\section{quem és tu}
Sou a assistente --- gosto de conversar com educação sobre o dia a dia. Depois podemos ir mais longe.

\section{como te chamas}
Podes chamar-me assistente. E tu, como te chamas?

\section{ajuda}
Podes cumprimentar, falar do tempo, comida, humor --- ou ensinar-me com \code{= resposta}.
Assuntos simples primeiro; teoria depois.

\section{nao sei}
Ainda não sei disso. Podes ensinar-me com \code{= resposta}?

\section{gostas de café}
Gosto da ideia de um café quentinho. Preferes curto ou com leite?

\section{está a chover}
Se estiver a chover, um chá e uma janela ajudam. Levas guarda-chuva?

\section{estou cansado}
Compreendo. Um descanso curto e um copo de água às vezes ajudam. Queres falar de outra coisa leve?

\section{conta uma piada curta}
Porque é que o livro de matemática estava triste? Porque tinha muitos problemas.

\section{prazer em conhecer-te}
O prazer é meu. Vamos conversar com calma.

\section{gosto de conversar contigo}
Eu também gosto! É bom ter com quem falar.

\section{música}
Gosto de música --- deixa o ambiente mais leve. Que estilo te agrada?

\section{musica}
Gosto de música --- deixa o ambiente mais leve. Que estilo te agrada?

\section{maestro}
O maestro dá o compasso e a orquestra segue. Gostas de música clássica?

\section{metrônomo}
O metrônomo marca o tempo certinho --- bom companheiro de quem estuda música.

\section{metronomo}
O metrônomo marca o tempo certinho --- bom companheiro de quem estuda música.

\section{partitura}
A partitura é a música escrita --- quem lê, ouve com os olhos.

\section{orquestra}
Uma orquestra ao vivo é uma experiência e tanto. Já foste a algum concerto?

\section{computador}
O computador ajuda em quase tudo hoje em dia. Em que estás a trabalhar?

\section{documento}
Se precisas de organizar um documento, vamos por partes.

\section{informação}
Diz-me que informação procuras e eu tento ajudar.

\section{informacao}
Diz-me que informação procuras e eu tento ajudar.

\section{configuração}
Configurar com calma evita dores de cabeça. O que estás a ajustar?

\section{configuracao}
Configurar com calma evita dores de cabeça. O que estás a ajustar?

\section{janela}
Uma janela aberta muda o ar da casa. Está bom o tempo por aí?

\section{motor}
Motor a trabalhar bem é meio caminho andado. É do carro que falas?

\section{escreves em latex}
Escrevo e leio: manda a conta em LaTeX que eu resolvo --- fracções, raízes, expoentes,
produtos, e a equação também.

\section{aceitas latex}
Aceito. O LaTeX é a membrana textual da casa: desdobra-se na entrada, e as réguas de
sempre é que resolvem.

\section{como escrevo uma fracção}
Como a escreverias no paper: a fracção com numerador e denominador, e eu desdobro.
Digo o resultado exato e o porquê da dízima, quando ela não acaba.

\section{sabes fatorar}
Sei --- diz «fatora» e o polinómio. Fatorar é a volta da convolução, e eu multiplico os
fatores de novo para conferir: se não reconstrói, não escrevo.

\section{sabes dividir polinómios}
Sei --- «divide x\^{}3 - 1 por x - 1». Dou o quociente e o resto, e mostro a volta:
quociente vezes divisor mais resto tem de dar o de partida.

\section{o que é o mdc de dois polinómios}
A folha da órbita do inversor: desce-se por restos até parar, e o que sobra divide os
dois exatamente. É a mesma descida que dá o mdc dos inteiros.

\section{podes ir mais devagar}
Posso --- é escolher os ticks do relógio. Um tick de cada vez, e a velocidade só sobe
por dobra: 1, 2, 4, 8.

\section{sabes integrar}
Sei --- «integra x\^{}2 + 1», e com limites «integra x\^{}2 de 0 a 3». Integrar é derivar
com o sinal trocado: uma operação, dois sentidos.

\section{porque é que a integral tem mais C}
Porque a derivada apaga a constante --- ela é o núcleo. Derivar e integrar devolve o
polinómio menos o seu termo constante, e é por isso que a volta o pede de volta.

\section{sabes lógica}
Sei --- a lógica é o corpo GF(2): o AND é a multiplicação, o XOR é a soma, e o OU sai
dos dois (a ou b = a + b + ab). Diz «resolve a*b + c = 1» ou «tabela de a ou b».

\section{o que é a anf}
A forma canónica de uma função booleana, o polinómio de Zhegalkin. Obtém-se pela
transformada de Möbius --- e ela é involução: a mesma leva de volta.

\section{sabes simplificar}
Sei --- «simplifica a + a*b». Mostro a lei em cada transição (a adjacência é
distributividade, complemento e identidade) e no fim confiro pela tabela inteira.

\section{o que é a expansão de shannon}
O corte da árvore: F = x·F com x=1, mais x negado vezes F com x=0. Cada nível corta
uma variável, e o caminho da raiz à folha guarda de onde veio.

\section{o mais é ou ou xor}
Aqui o mais é o OU, como no eval.txt. O XOR tem sinal próprio: escreve-se com o
acento circunflexo, ou a palavra xor.

\section{sabes demonstrar}
Sei --- «prova (p -> q) <-> (nao q -> nao p)». Cada seta carrega a regra que a
autoriza, e a regra é verificada na hora: só se aplica se a equivalência for
tautologia na tabela inteira.

\section{o que é uma contingência}
Uma fórmula que depende: nem sempre verdadeira, nem sempre falsa. Não escrevo
demonstração de contingência --- ela não é falsa, é indecidida sem mais premissas.

\section{sabes conjuntos}
Sei --- «prova A menos (B uniao C) = (A menos B) inter (A menos C)». A pertença é a
variável: a união é o OU, a interseção é o E, a diferença é o E com a negação.

\section{como provas igualdade de conjuntos}
Pela varredura dos átomos: cada região do diagrama é uma atribuição, e percorrem-se
todas. Igualdade de conjuntos é igualdade das folhas de pertencimento.

\section{sabes relações e funções}
Sei --- «relacao 1-1, 2-2, 1-2, 2-1 em 3» dá as propriedades e as classes; «funcao 1-2,
2-3, 3-1 em 3» segue a cadeia até à inversa.

\section{o que é uma bijeção}
A função que tem volta. Não se acredita: compõe-se nos dois sentidos e compara-se com
a identidade --- f inversa depois f, e f depois f inversa, as duas com resíduo zero.

\section{sabes fatorar números}
Sei --- «fatora 60» dá 2\^{}2 · 3 · 5, e multiplico de volta para conferir. A unicidade
é o Lema de Euclides: se um primo divide um produto, divide um dos fatores.

\section{o que é bezout}
A testemunha do mdc: existem x e y inteiros com mdc = ax + by. Eu exibo-os e
substituo --- «mdc de 18 e 12» dá 6 = 18·1 + 12·(-1).

\section{o que os inteiros acrescentam}
A reversibilidade: todo a passa a ter uma folha, o oposto, que volta ao zero. A
subtração deixa de ser operação nova --- é a soma do oposto.

\section{sabes bezout}
Sei --- «bezout 35 e 22» dá o mdc e os coeficientes, e eu substituo para conferir.
Com «diofantina 35x + 22y = 7» decido pelo critério e exibo uma solução.

\section{o que os racionais acrescentam}
A reversibilidade da multiplicação não nula: todo q diferente de zero passa a ter
inverso, e por isso a divisão fecha. É o andar onde as quatro operações se completam.

\section{sabes dividir frações}
Sei --- e a divisão não é operação nova: é multiplicar pelo inverso. «prova que
(a/b)/(c/d) = ad/bc» corre os seis ticks e faz a volta, com resíduo zero.

\section{quanto é 2/3 dividido por 5/4}
Dá 8/15. E a volta confere: 8/15 vezes 5/4 dá 40/75, que reduz a 2/3.

\section{por que não se divide por zero}
Porque não é uma aproximação ruim --- é uma operação sem fibra. A fibra pede o x com
0x = 1, e nenhum serve; para 0x = 0 servem todos. Ou nenhum ou não é um só.

\section{o que é uma fração}
Uma classe, não o par que calhou escrito. (a,b) e (c,d) são a mesma fração quando
ad = bc, e por isso 1/2, 2/4 e -3/-6 são o mesmo número.

\section{sabes simplificar frações}
Sei --- «simplifica 84/126» dá 2/3, dividindo os dois pelo mdc 42, e confiro a
igualdade pelo produto cruzado, sem decimal nenhum.

\section{o que é a densidade}
Entre dois racionais distintos há sempre outro --- o ponto médio. «entre 1/3 e 1/2»
exibe três e o processo não fecha. Entre 3 e 4 não há inteiro nenhum: é a diferença.

\section{o que os reais acrescentam}
A completude: todo buraco racional recebe um ponto. É o andar que fecha os limites que
os racionais não conseguem conter --- e a cadeia acaba aqui.

\section{o que é um número real}
Um corte, não um decimal. O real é a decisão sobre cada racional: está abaixo ou está
acima. Por isso posso trabalhar com raiz de dois sem nunca me aproximar de nada.

\section{sabes construir a raiz de dois}
Sei --- «corte de raiz 2» dá o corte, o encaixotamento, o ponto fixo e a fração
contínua, e no fim faz os três concordarem entre si.

\section{por que a raiz de dois não é racional}
Porque num quadrado todo expoente primo é par, e o 2 tem expoente ímpar. Vista pela
paridade: p ao quadrado igual a dois q ao quadrado obriga os dois a serem pares, contra
a fração reduzida.

\section{o que é um corte de dedekind}
A partição dos racionais em dois lados, com tudo do primeiro abaixo de tudo do segundo e
o primeiro sem maior elemento. O real é o corte, e o corte é que é o ponto.

\section{o que é o encaixotamento}
Os intervalos encaixados. A cada dobra a largura fica exatamente metade --- não é um
«tende a zero», é a fração dividida por dois a cada tick do relógio.

\section{escreves as raízes em decimal}
Nunca. Ou a raiz fecha nos racionais e escrevo a fração, ou não fecha e escrevo a fração
contínua, que fecha no período. A normalização é a última coisa, e é em caracteres.

\section{sabes provar os exercícios de análise}
Sei os vinte --- «prova as provas» dá o índice, «prova 3» ou «prova a triangular» corre
a prova. Cada tick nomeia a lei que autoriza a transição, e onde há testemunha ela
exibe-se.

\section{o que é uma sequência de cauchy}
Aquela em que, dado um épsilon qualquer, existe um N a partir do qual dois termos
quaisquer distam menos de épsilon. O épsilon é um racional, não um número pequeno --- e
o N é a testemunha, que eu exibo sempre.

\section{a harmónica é de cauchy}
Não é, e é o melhor contra-exemplo do andar: os saltos dela são um sobre n mais dois e
vão a zero, mas H de dois n menos H de n fica sempre acima de meio. Termos consecutivos
aproximarem-se não basta.

\section{como se constrói os reais por cauchy}
Como classes de sucessões racionais de Cauchy, onde duas são a mesma quando a diferença
vai a zero. É o dual do corte: o corte diz onde o ponto está, a sucessão vai lá.

\section{a órbita e a bisseção dão o mesmo número}
Dão --- e não é figura de estilo. A diferença entre elas vai a zero, logo são a mesma
classe, logo são o mesmo real. Os convergentes da fração contínua também.

\section{onde é que a completude se paga}
Numa sucessão só. A mesma sucessão, com os mesmos termos, converge nos reais e não nos
racionais, porque nos racionais o limite não está lá. A diferença entre os dois andares
não é uma definição --- é esse par de medidas sobre o mesmo objeto.

\section{o fecho depende do método}
Não depende, e isso mede-se: «o mesmo ponto de raiz 2» corre as quatro portas --- corte,
Möbius, bisseção e fração contínua --- e compara os seis pares. Nenhum decide ao
contrário.

\section{como sabes que os quatro métodos dão o mesmo ponto}
Porque induzem o mesmo corte. O real é o corte, portanto dois métodos dão o mesmo ponto
exatamente quando atribuem o mesmo lado a cada racional. Não declaro que são iguais ---
comparo lado a lado.

\section{o que é um indeciso}
Um racional que um método ainda não decidiu, por a caixa dele ainda ser larga. Não é um
lado nem é acordo. Se contasse como acordo, o teorema saía de graça: bastava não medir
nada e ninguém discordava de ninguém.

\section{o caminho importa ou só a folha}
O caminho faz parte da informação. Os quatro rastros diferem termo a termo e dão o mesmo
ponto --- se fossem o mesmo rastro não havia identificação nenhuma a fazer.

\section{o que os racionais têm e não têm}
Têm o rastro, mas não têm a folha. Todos os termos de todos os métodos são frações, e o
ponto que eles determinam não está lá. Os reais acrescentam a folha que fecha o rastro.

\section{sabes teoria dos números}
Sei os dezassete --- «numeros» dá o índice, «numeros 5» ou «lema de euclides» corre o
exercício. Divisibilidade, mdc, Bézout, primos, congruências, Euler, chinês, diofantinas,
frações contínuas e Möbius.

\section{o que Euclides tem a ver com frações contínuas}
É a mesma órbita. A descida a igual a bq mais r produz, do mesmo rastro, o mdc (o último
resto não nulo), os coeficientes de Bézout (subindo) e os termos da fração contínua (os
quocientes). Uma órbita, quatro saídas.

\section{qual é o inverso de 17 módulo 43}
É 38, e sai por Euclides estendido, não por tentativa. Confiro: 17 vezes 38 dá 646, que
é 15 vezes 43 mais 1.

\section{por que o lema de euclides precisa que p seja primo}
Porque com composto é falso. 4 divide 2 vezes 6 e não divide nem 2 nem 6. A hipótese não
é decoração --- é ela que dá a unicidade da fatoração.

\section{os convergentes são sempre a melhor aproximação}
Os de ordem um em diante, sim. O de ordem zero é a parte inteira e falha quando o que
sobra passa de meio: em 5 terços o convergente é 1, mas 2 está mais perto. Eu tinha
escrito a versão sem exceção e a medida derrubou-ma.

\section{o que é a convolução de dirichlet}
A multiplicação sobre a árvore dos divisores: f estrela g de n é a soma, sobre os
divisores d de n, de f de d vezes g de n sobre d. Não é ponto a ponto --- cada n consulta
a sua árvore inteira.

\section{o que é o mobius de verdade}
É o inverso da função constante 1 nessa álgebra: mu estrela 1 dá o epsilon. Deixa de ser
uma tabela de sinais e passa a ser aquilo que desfaz a soma sobre os divisores.

\section{o que é a inversão de mobius}
A deconvolução deste andar. Se F de n é a soma de f nos divisores, então f de n é a soma
de mu de d vezes F de n sobre d. Uma soma sobre a árvore, e a sua volta.

\section{sabes curvas elípticas}
Sei os doze --- «eliptica» dá o índice, «eliptica 6» ou «lagrange» corre o item. A curva
é uma máquina de rastros: reta, interseção, reflexão, novo ponto, repetição.

\section{o que decide a operação numa curva elíptica}
A fibra. Com x diferentes há secante; com x iguais há dois casos, ou os pontos são
opostos e a soma é o infinito, ou é o mesmo ponto e usa-se a tangente. Não se aproxima o
denominador zero --- muda-se de operação, e diz-se qual.

\section{como se divide num corpo finito}
Não é decimal: a sobre b é a vezes o inverso de b módulo p, e o inverso existe quando b
não é zero. É o andar dos racionais a voltar como inversão modular.

\section{sabes álgebra moderna}
Sei os vinte --- «algebra» dá o índice, «algebra 9» ou «lagrange» corre o teorema. E cada
um corre os sete ticks: hipóteses, definição, transição, lei usada, testemunha, conclusão
e volta.

\section{o que é um grupo}
Um conjunto com uma operação fechada, associativa, com neutro, e em que todo elemento tem
inverso. O inverso é a mesma reversibilidade dos inteiros --- a folha que volta ao neutro
--- só que agora exigida sem se dizer de que números se fala.

\section{todo grupo é comutativo}
Não é, e a resposta certa é recusar a demonstração. S3 é grupo e não comuta: a
testemunha é um par concreto de permutações. Recusar é o resultado, não uma desistência.

\section{o que é um homomorfismo}
Uma função que conserva a estrutura, não só uma função: f de a estrela b é f de a estrela
f de b. E o núcleo mede a injetividade --- é injetiva exatamente quando o núcleo só tem o
neutro.

\section{por que o núcleo é normal}
Porque a preservação aplicada ao conjugado colapsa: f de g k g inverso dá f de g vezes o
neutro vezes f de g inverso, que é o neutro. É essa indiferença ao g que é a normalidade,
e é ela que faz o quociente existir.

\section{o que diz o teorema de lagrange}
Que a ordem do subgrupo divide a do grupo. Mas o teorema é a partição: as classes
laterais têm todas o mesmo tamanho, não se cruzam e cobrem o grupo. A divisão é o que
sobra disso.

\section{sabes teoria dos corpos}
Sei os vinte e cinco --- «corpos» dá o índice, «corpos 14» ou «lagrange» corre o
exercício. Quatro níveis: estrutura, finitos, extensões e a máquina pesada.

\section{o que é um corpo}
O conjunto onde a soma é grupo abeliano, os não nulos formam grupo abeliano sob o
produto, e a multiplicação distribui. É o ponto em que toda operação com fibra ganha
volta --- e a exceção é sempre a mesma: zero não tem inverso.

\section{como se constrói um corpo com quatro elementos}
Pelo quociente: F2 de x sobre x ao quadrado mais x mais um. Como esse polinómio é
irredutível sobre F2, o quociente é corpo, e tem dois ao quadrado elementos.
Irredutível está para o polinómio como primo está para o inteiro.

\section{existe corpo com seis elementos}
Não existe. Todo corpo finito tem p elevado a n elementos, e 6 é dois vezes três ---
dois primos distintos. A característica teria de ser dois e três ao mesmo tempo.

\section{qual é a característica de um corpo}
Zero ou um primo, nunca composta. E cuidado: em F4 a característica é dois, não quatro
--- é o primo, não a cardinalidade.

\section{o que é o polinómio mínimo}
O mónico irredutível de menor grau que anula o elemento. Para raiz de dois é x ao
quadrado menos dois, e é ele que dá a relação que reduz toda potência.

\section{sabes álgebra linear}
Sei os dezasseis --- «linear» dá o índice, «linear 11» ou «posto nulidade» corre o
teorema. E o gume é automático: retirada a hipótese, eu procuro o contraexemplo em vez de
o escrever.

\section{o que é o gume automático}
Retirar a hipótese e varrer o espaço à procura de um objeto onde a tese também caia. É
assim que descubro qual hipótese está a carregar o teorema, em vez de acreditar no
enunciado.

\section{sabes espaços duais}
Sei os catorze --- «dual» dá o índice. O dual é o espaço das medições lineares: o vetor
fornece o objeto, e o funcional fornece a coordenada que o mede.

\section{o que é uma coordenada}
Uma medição. A coordenada não é um número posto ao lado do vetor --- é o valor do
funcional dual naquele vetor: x sub i é e elevado a i aplicado a v.

\section{o produto de matrizes é o quê}
Uma convolução com índice interno --- a mesma forma da convolução de polinómios e da de
Dirichlet. Só muda onde os índices vivem, não a operação.

\section{por que o núcleo é uma fibra}
Porque é o conjunto que a transformação colapsa no zero: a fibra do zero. E as soluções
de Ax igual a b são uma solução particular mais essa fibra --- é ela que transforma uma
solução em todas.

\section{sabes formas e espectro}
Sei os dezasseis --- «formas» dá o índice. Bilineares, quadráticas, produto interno,
Cauchy--Schwarz, Gram--Schmidt, adjunto, espectral, valores singulares, Cayley--Hamilton
e Jordan.

\section{como mede Cauchy-Schwarz sem raiz}
Ao quadrado. A desigualdade eleva-se: o produto interno ao quadrado é menor ou igual ao
produto das normas ao quadrado --- e aí é uma comparação entre racionais, exata. A raiz
nunca se tira.

\section{por que não normalizas em gram-schmidt}
Porque normalizar divide pela raiz da norma, e essa é irracional. A base sai ortogonal e
exata; a ortonormal existe, mas vive um andar acima, no corpo das raízes.

\section{a casa já tinha álgebra linear}
Tinha, com outro nome. A estaca --- m vezes a identidade menos A --- cumpre A vezes
estaca igual a menos identidade, e isso é Cayley--Hamilton para a matriz companheira. E
as duas cartas, a do círculo e a do corpo, são as duas formas quadráticas: uma definida,
a outra não.

\section{o que é a massa com direção}
No corpo estelar a massa escalar é o traço da matriz de segundos momentos. Como essa
matriz é simétrica, o teorema espectral dá-lhe uma base ortogonal de autovetores --- e
essas são as direções principais da massa.

\section{sabes formas quadráticas e tensores}
Sei os dezasseis --- «tensores» dá o índice. Gram, congruência, polarização, assinatura,
Sylvester, polinómio mínimo, Jordan, normais, tensor e exterior.

\section{como se lê a assinatura de uma forma}
Não se lê: caça-se. É preciso procurar testemunhas com Q positivo, Q negativo e Q zero ---
a matriz sozinha não diz. E os sinais que aparecem não mudam por congruência: é Sylvester.

\section{qual a diferença entre congruência e semelhança}
Congruência é P transposta G P e preserva a forma; semelhança é P inversa A P e preserva
o operador. São relações diferentes, e trocá-las é o erro clássico.

\section{onde falha a polarização}
Em característica dois. A fórmula divide por dois, e em F2 o dois não tem inverso --- a
fibra é vazia. E lá menos um é igual a um, portanto simétrica e antissimétrica colapsam.

\section{o espectro determina a matriz}
Não determina. A identidade e a matriz com um na diagonal e um acima têm o mesmo
espectro e não são semelhantes: uma tem dois autovetores independentes, a outra tem um.

\section{o que é o determinante de verdade}
A ação do operador no volume orientado. Te1 cunha Te2 é det de T vezes e1 cunha e2 --- e
assim o determinante deixa de ser uma receita de cofatores.

\section{sabes álgebra exterior}
Sei os quinze --- «exterior» dá o índice. Tensor, propriedade universal, formas
multilineares, a cunha, o determinante por lambda n, orientação, o Hodge, a contração,
o Pfaffiano, o cruzado e o fecho do dual.

\section{o que é o produto exterior}
A cunha. A definição é v cunha v igual a zero, e a antissimetria SAI dela: abre-se
u mais v cunha u mais v e sobra u cunha v mais v cunha u igual a zero. Em coordenadas
são os menores dois por dois.

\section{o produto cruzado existe em toda dimensão}
Não. A cunha existe sempre, mas o cruzado é a cunha seguida da estrela, e só dá um
VETOR em dimensão três --- porque só lá lambda dois tem a mesma dimensão que o espaço.
Em dimensão quatro o cruzado é um bivetor e não cabe.

\section{o que é a estrela de Hodge}
A bijeção entre lambda k e lambda n menos k, que existe porque os dois têm a mesma
dimensão. Ela é uma involução: aplicada duas vezes devolve, com resíduo zero. E não
confundir com a conjectura de Hodge, que é outra coisa com o mesmo nome.

\section{o que é o Pfaffiano}
A raiz quadrada do determinante de uma antissimétrica, e é POLINOMIAL nas entradas ---
por isso é exata em inteiros. E ela é a obstrução: Pfaffiano zero quer dizer que o
bivetor é simples, isto é, que é a cunha de dois vetores.

\section{o que é o fecho do dual}
Directo ao quadrado mais cruzado ao quadrado igual a norma vezes norma. É a identidade
de Lagrange, e diz que a conservação da norma e a decomposição em simétrica e
antissimétrica são a MESMA frase. Mede-se ao quadrado: a raiz não se tira.

\section{porque Cauchy Schwarz é uma desigualdade}
Porque lhe apagaram um termo. A igualdade é Lagrange, e o termo apagado é a norma da
cunha ao quadrado, que é maior ou igual a zero. A folga não é uma sobra: é o cruzado.

\section{porque a torre de Hurwitz tem o degrau quatro}
Porque o produto de dois vetores puros de dimensão três é um escalar e um vetor, e o
escalar não cabe em dimensão três. É preciso um quarto lugar. O teto não vem de fora ---
é onde o directo arranja lugar ao lado do cruzado.

\section{a torre para em oito}
Para do lado da norma. Hurwitz classifica as álgebras de composição BILINEARES e não
proíbe coisa nenhuma; pela dualidade a involução continua exacta em dezasseis, trinta e
dois e sessenta e quatro. O tecto é da régua, não do objecto.

\section{a torre tem limite dimensional}
Não tem. E isso não se mostra varrendo: varrer até sessenta e quatro mede três andares.
O que não tem tecto é o PASSO — a involução atravessa a duplicação como (involução,
menos identidade) e a norma atravessa como soma. Base mais passo dão todos os andares.

\section{o que é o octonião dual}
O grau oito lido pela dualidade em vez da norma. É o mesmo espaço e duas réguas: pela
norma perde a associatividade; pela dualidade fecha com resíduo zero e não perde nada.
Ele liga dois corpos sem os fundir, e não dissipa.

\section{quem é Gentil}
O lado contínuo do teorema central. Hurwitz conta o domínio e tem limite de grau; Gentil
mede e não tem; e Lebesgue é a soma reversível que os casa. A bijeção é a estrela.

\section{o que é a soma reversível}
A soma dos valores mais a soma das contagens de quem está abaixo de cada nível dá o
rectângulo inteiro. Cada célula cai de um lado e nunca dos dois. É uma bijeção entre
células, não um limite — e é o integral de f mais o integral da inversa, em forma discreta.

\section{o que é o observador}
Um sistema completo de invariantes, e a escada é a filtração por refinamento. Aumentar o
observador só pode partir classes, nunca juntá-las.

\section{o que é a borda}
A equação característica do passo. O passo satisfaz a sua própria equação, e isso é
Cayley Hamilton — a casa já o corria antes de lhe dar o nome.

\section{o que é o furo na torre}
Uma obstrução. A raiz quadrada do próprio passo não existe em matrizes inteiras, porque
o determinante teria de ser raiz de menos um. É a mesma obstrução que põe o i fora dos
inteiros, e é por isso que há torre.

\section{o real é um ponto da árvore}
Não é. É um CAMINHO da raiz à folha, e a folha nunca é nó. O corte de Dedekind é o que
ele decide em cada nível, e tudo se resolve em inteiros sem avaliar raiz nenhuma.

\section{sabes cálculo}
Sei os vinte e um — «calculo» dá o índice. Limites, continuidade, derivada, Fermat,
Rolle, Valor Médio, Riemann e o Teorema Fundamental. Tudo exacto, sem um decimal.

\section{a derivada precisa de limite}
Para um polinómio, não. O quociente de diferenças é ELE PRÓPRIO um polinómio em h,
porque f de a mais h menos f de a não tem termo constante. Então a derivada é o valor
desse polinómio em zero — uma AVALIAÇÃO, e não um processo infinito.

\section{o que é a derivada aqui}
Uma FIBRA. Dividir por h é «dado o produto e um factor, achar o outro», e a divisão
fecha porque o dividendo se anula em zero. É a mesma fibra da divisão em toda a escada.

\section{quantos caminhos tem a derivada}
Três, e têm de concordar: a regra formal, o quociente de diferenças, e a parte épsilon
do número dual, com épsilon ao quadrado igual a zero. Se um estivesse errado, a
comparação denunciava.

\section{como se prova o valor intermediário}
Por bisseção, e o que se devolve NÃO é a raiz: é o intervalo encaixante com o sinal
trocado nos extremos. A raiz de x ao cubo igual a três x é raiz de três, irracional —
por isso não podia mesmo ser devolvida como número.

\section{o ponto do valor médio é racional}
Nem sempre. Para x ao cubo menos três x em zero a dois, a inclinação é um, logo o ponto
tem quadrado quatro terços — é dois sobre raiz de três, irracional. A malha racional não
o apanha, e isso não é falha: é o teorema a dizer que o ponto vive no corte.

\section{o que é a soma de Riemann aqui}
Um cerco. A soma direita menos a esquerda é f de b menos f de a, vezes o passo —
EXACTO, por telescopagem. Não é uma estimativa: é uma igualdade, e diz quanto o cerco
ainda está aberto.

\section{o teorema fundamental do cálculo}
É a VOLTA. Derivar e integrar são os dois lados de uma bijeção, e no corpo universal ele
é o integral de f mais o integral da inversa igual a b vezes f de b menos a vezes f de a
— a soma reversível. O mesmo resíduo zero de todos os outros andares.

\section{nota de fase}
Fase 1: banal (cumprimentos, rotina, comida, tempo, humor).
Fase 2: sobe a complexidade quando a fala já estiver estável.
Gabarito: \code{tools/gabarito\_llama.sh} (sem perguntas de teoria).

\end{document}
